\chapter{Signs, duals and coefficient completions}
\label{mht:ch:signs-duals}

A coefficient calculation enters a field theory through a map of
complexes. Its grading determines the signs of that map, and its
topology determines which inputs the map accepts. We give the
conventions needed for the Hamiltonian and boundary constructions.
The examples distinguish a coordinate change from a change of
representation or completion.

\section{Degrees and deformation parameters}

For a cochain complex, use
\[
 (V[r])^j=V^{j+r},\qquad
 d_{V[r]}=(-1)^r d_V.
\]
Thus an element of degree zero in $V$ has degree $-1$ in $V[1]$.
The tensor differential and interchange law are
\[
 d(v\otimes w)=dv\otimes w+(-1)^{|v|}v\otimes dw,
 \qquad
 v\otimes w\longmapsto(-1)^{|v||w|}w\otimes v.
\]
Parity is cohomological degree modulo two. When a complex is written
homologically, its differential lowers degree. We state this change
before introducing its generators.

The matrix relation variable $\Psi$ has cohomological degree $-1$,
and the gauge ghost $c$ has degree $1$. Both are odd. Form degree
on $\mathbb R^2\times\mathbb C^2$ is an additional grading before
the field complex is totalized. In a tensor product of form and
coefficient complexes, moving a coefficient past a form contributes
the sign from their declared degrees.

Several parameters occur in the comparison. Write $t$ for the Rees
parameter of a filtered operator algebra and $h$ for the coefficient
in a Weyl relation. The loop parameter in the quantum master
equation is $\hbar_Q$.

An equivariant parameter is specified by its
cohomological degree and vector field. The polar constant $\kappa$
is the value of a chosen residue normalization. These parameters
can be related by a stated scalar map. The grading convention
does not identify them.

\section{Cyclic derivatives and the relation differential}

Let $F=\mathbb C\langle x,y,c,\psi\rangle$, with degrees
$0,0,1,-1$, respectively. Its graded cyclic quotient is the vector
space
\[
 F_{\mathrm{cyc}}=F/[F,F]_{\mathrm{gr}},\qquad
 [uv]=(-1)^{|u||v|}[vu].
\]
The denominator is the linear span of graded commutators. It is not
the ideal imposing graded commutativity on the free algebra.

For a word $a_1\cdots a_n$, define its left cyclic derivative by
\[
 \partial_q[a_1\cdots a_n]
 =\sum_{a_i=q}
 (-1)^{|a_1\cdots a_{i-1}|(|q|+|a_{i+1}\cdots a_n|)}
 a_{i+1}\cdots a_na_1\cdots a_{i-1}.
\]
This sign moves the preceding block past the selected factor and
the following block. Fix the even graded-skew pairing
\[
 \omega(x,y)=1,\quad\omega(y,x)=-1,
 \qquad \omega(c,\psi)=\omega(\psi,c)=1.
\]
With $\iota_{X_G}\omega=-dG$, the induced necklace bracket is
\[
 \begin{split}
 \{G,H\}_{\mathrm{neck}}=
 \bigl[&\partial_xG\,\partial_yH-\partial_yG\,\partial_xH\\
 &+(-1)^{|G|}(\partial_cG\,\partial_\psi H
                         +\partial_\psi G\,\partial_cH)\bigr].
 \end{split}
\]
In particular $\{[x],[y]\}=[1]$ and
$\{[c],[\psi]\}=-[1]$. The symmetry of the pairing on the two
odd factors is required by graded skew symmetry.

The full matrix BRST differential of
\eqref{mht:eq:matrix-classical-brst} uses graded matrix products:
\[
 QX=[X,c],\qquad QY=[Y,c],\qquad Qc=-c^2,
 \qquad Q\Psi=[X,Y]-[c,\Psi].
\]
The relation term is equivariant for conjugation. For example,
putting $M=[X,Y]$ gives $QM=[M,c]$, and
\[
 Q[c,\Psi]=[Qc,\Psi]-[c,Q\Psi].
\]
Since $[c,[c,\Psi]]=[c^2,\Psi]$, substitution gives $Q^2\Psi=0$.
The same computation for $X,Y$, and $Q(c^2)=(Qc)c-c(Qc)$,
gives $Q^2=0$ on all generators. The odd Leibniz rule extends this
identity to the polynomial algebra.

Without ghosts, set $dx=dy=0$ and $d\psi=xy-yx$. For even words
$u,v$,
\begin{equation}\label{mht:eq:sign-two-psi}
 d[\psi u\psi v]
   =[\psi vxyu]-[\psi vyxu]-[\psi uxyv]+[\psi uyxv].
\end{equation}
The second replacement crosses one odd factor and has a minus
sign. The first replacement can be rotated past the remaining
$\psi$ without a further sign because the moved block is even.
Thus interchanging $u,v$ negates the result, as required by graded
cyclicity.

\section{Suspension and the ordered bar}

Let $A$ be an augmented differential graded algebra in cohomological
degrees, and let $\bar A$ be its augmentation ideal. Write
$sa\in\bar A[1]$, so $|sa|=|a|-1$. The tensor coalgebra
\[
 BA=\bigoplus_{n\geq0}(\bar A[1])^{\otimes n}
\]
has deconcatenation coproduct, including the empty word. Its
differential is the coderivation determined by
\[
 b_1(sa)=-s(da),\qquad
 b_2(sa\otimes sb)=(-1)^{|a|}s(ab).
\]
Inserting a degree-one map after $sa_1,\ldots,sa_{i-1}$ contributes
$(-1)^{\sum_{j<i}(|a_j|-1)}$.

\begin{proposition}\label{mht:prop:suspended-bar-signs}
These formulas give $b^2=0$. For degree-zero closed elements,
\[
 b[a|b|c]=[ab|c]-[a|bc].
\]
Consequently its second boundary is $[(ab)c-a(bc)]=0$.
\end{proposition}

\begin{proof}
On one input the assertion is $d^2=0$. On two inputs the terms
containing $d(ab)$ cancel the two terms containing $da$ and $db$,
using $d(ab)=(da)b+(-1)^{|a|}a\,db$.
On three inputs the two multiplication terms have opposite signs
and cancel by associativity. On disjoint input blocks the two
degree-one insertions anticommute. These cases exhaust the square
of the coderivation at every word length.
\end{proof}

The empty word is the coalgebra coaugmentation. An algebra
augmentation is a separate map $A\to\mathbb C$ used to define
$\bar A$. A local vacuum or a prepared state in a field theory
requires its own map. Neither is selected by the empty bar word.

\section{Chevalley coordinates and cotangent shifts}

Let $\mathfrak l$ be a finite-dimensional Lie algebra with
$[e_I,e_J]=C^K_{IJ}e_K$. In
$\mathfrak l\ltimes\mathfrak l^\vee[1]$, the dual summand has
degree $-1$ and is abelian. Its Chevalley algebra has coordinates
\[
 |c^I|=1,\qquad |u_I|=2,
\]
and differential
\[
 dc^K=-\tfrac12 C^K_{IJ}c^Ic^J,
 \qquad du_K=C^L_{JK}u_Lc^J.
\]
Give $u_I$ weight one and $c^I$ weight zero. Subtracting twice this
weight from cohomological degree makes the $u_I$ even coordinates
of degree zero, without changing the degree of the differential.

For the alternative shift
$\mathfrak l\ltimes\mathfrak l^\vee[-1]$, these degree-zero
coordinates occur without regrading. The two original graded Lie
algebras have different cotangent degrees. Their coordinate
comparison is the stated regrading, not a comparison of their
graded field pairings.

Write $O_I$ for linear functions on $\mathfrak l^\vee$ and
$\partial^I$ for constant vector fields, with
$[\partial^I,O_J]_S=\delta^I_J$. The linear Poisson tensor is
\[
 \pi=\tfrac12 C^K_{IJ}O_K\partial^I\partial^J.
\]
Its differential $d_\pi=[\pi,-]_S$ satisfies
\[
 d_\pi\partial^K=-\tfrac12C^K_{IJ}\partial^I\partial^J,
 \qquad d_\pi O_K=C^L_{JK}O_L\partial^J.
\]
Thus $c^I\mapsto\partial^I$, $u_I\mapsto O_I$ intertwines the
two differentials on generators and hence on their symmetric
algebras. The vector field $d_\pi O_I$ is generally not
$\partial^I$.

For $[e_1,e_2]=e_2$, the equations become
\[
 dc^1=0,\quad dc^2=-c^1c^2,
 \qquad du_1=-u_2c^2,\quad du_2=u_2c^1.
\]
Then $d^2u_1=-u_2c^1c^2+u_2c^1c^2=0$.
This two-dimensional example detects a reversal of the coadjoint
sign. For an abelian Lie algebra every $d_\pi O_I$ vanishes, while
the constant fields remain nonzero.

\section{All antifield degrees in the cyclic-word complex}

Use homological degrees in this section. Let
$R=\mathbb C\langle x,y,\psi\rangle$, with $|\psi|=1$ and
$d\psi=xy-yx$. Give the factors weights $(1,0),(0,1),(1,1)$.
In bidegree $(r,s)$, a degree-$p$ cyclic word has $p$ occurrences
of $\psi$, $r-p$ of $x$, and $s-p$ of $y$. Denote this finite
complex by $K_{r,s}$.

There is an elementary model for all of its degrees. For $s>0$,
put
\[
 B_s=\mathbb C[v_0,\ldots,v_{s-1}]
       \otimes\Lambda(e_0,\ldots,e_{s-1}),
 \qquad de_i=v_i-v_{i+1},\quad dv_i=0,
\]
with indices modulo $s$. Give each generator weight one, each
$v_i$ homological degree zero and each $e_i$ degree one. A
monomial of weight $r$ records a labeled cyclic word. The factor
$v_i^{a_i}$ records $a_i$ copies of $x$ immediately before the
$i$th marker. This marker is $\psi$ when $e_i$ occurs and $y$ otherwise.

Replacing the marker $\psi$ by $xy-yx$ increases the multiplicity
at vertex $i$ or vertex $i+1$, with opposite signs. Ordering the
occupied markers gives exactly the exterior signs of $B_s$.
Rotation of the $s$ labels gives the graded cyclic relation,
because interchanging blocks with $a,b$ occupied markers changes
orientation by $(-1)^{ab}$. Hence
\[
 K_{r,s}\cong (B_s^{\mathrm{wt}=r})_{\mathbb Z/s\mathbb Z}.
\]

\begin{proposition}\label{mht:prop:cyclic-koszul-circle}
For $r,s>0$, the homology of $K_{r,s}$ is one-dimensional in
degrees zero and one, and zero in every higher degree. A
degree-one generator is
\[
 \Psi_{r-1,s-1}
   =\sum_{w\in\mathsf W_{r-1,s-1}}[\psi w],
\]
where $\mathsf W_{a,b}$ consists of the words with $a$ factors
$x$ and $b$ factors $y$.
\end{proposition}

\begin{proof}
Set $z=v_0$, $a_i=v_i-v_{i+1}$ for $0\leq i<s-1$, and
$\eta=\sum_i e_i$. These are invertible linear changes of
the even and odd variables. They give
\[
 B_s\cong\mathbb C[z]\otimes\Lambda(\eta)\otimes
 \bigotimes_{i=0}^{s-2}
       (\mathbb C[a_i]\otimes\Lambda(e_i),\ de_i=a_i).
\]
The last factors are one-variable Koszul resolutions. Their
positive-weight parts contract by polynomial division by $a_i$.
Thus $H_*(B_s)=\mathbb C[z]\otimes\Lambda(\eta)$.

At weight $r>0$ the surviving classes are $z^r$ and
$z^{r-1}\eta$. Rotation fixes $\eta$ and sends every $v_i$ to
a representative of $[z]$, so it acts trivially on homology.
Coinvariants for the finite cyclic group are exact over
$\mathbb C$, proving the dimension assertion.

Every necklace has equally many $xy$ and $yx$ transitions, so
$d\Psi_{a,b}=0$. The functional taking $[\psi w]$ to $1$
kills every two-antifield boundary by
\eqref{mht:eq:sign-two-psi}. Its value on $\Psi_{a,b}$ is
$\binom{a+b}{a}$, which is nonzero. This identifies the generator.
\end{proof}

If $r=0$ or $s=0$, only the unique degree-zero necklace remains.
Summing all weights gives
\[
 H_*(R/[R,R]_{\mathrm{gr}})
   \cong\mathbb C[x,y]\oplus\mathbb C[x,y]\theta,
 \qquad |\theta|=1,\quad\mathrm{wt}(\theta)=(1,1).
\]
This is an identification of graded vector spaces. A multiplication
on the polynomial label space is additional structure on those
labels. The cyclic quotient itself is not an associative quotient
algebra.

The classes $[1]$ and $[\psi]$ occur in different bidegrees.
Removing the scalar augmentation removes $[1]$ and retains
$[\psi]$. Removing the latter requires another quotient. Nor is
the one-dimensional label in each bidegree the whole symmetric
algebra on the Hamiltonian Lie algebra: products of Hamiltonian
generators create additional summands. An identification with a
PBW component must specify that component and its map.

\section{Trace normalization and the central rank}

For an ungraded associative algebra $A$, use $N\times N$ matrices.
Write $E_{ij}a$ for the matrix with entry $a$ in position $(i,j)$. In ordinary homological CE
chains set
\[
 \Theta_1(a)=\sum_iE_{ii}a,
 \qquad \Theta_2(a,b)=\sum_{i,j}E_{ij}a\wedge E_{ji}b.
\]
The boundary is the matrix commutator, so
\[
 d\Theta_2(a,b)=N\Theta_1(ab-ba).
\]
Both products sum over a free matrix index. Thus the raw reverse
trace formula has a rank factor even in length two.

The primitive cyclic comparison in
Theorem~\ref{mht:thm:stable-cyclic-primitives} instead uses stable
scalar-matrix coinvariants and single cycles with distinct indices.

The empty cyclic word $\nu$ is central for the necklace bracket,
with $[x,y]=\nu$. It is a primitive generator before rank is fixed.
Matrix evaluation sends it to $N1$. Consequently, for
$X=\operatorname{Tr}x$ and $Y=\operatorname{Tr}y$,
\[
 \{X^2,Y\}=2NX.
\]

Projection that kills multitraces and scalar values kills $X^2$
but retains this bracket. The Lie projection in a symmetric
algebra on a zero-central-character Lie algebra must therefore be
taken before the fixed-rank specialization. These are different
reductions of the scalar direction.

\section{The residue dual of formal Hamiltonians}

Put $\mathcal O=\mathbb C[[u,v]]$ and
$\mathfrak m=(u,v)$. The formal Hamiltonian Lie algebra is
$\mathfrak h=\mathcal O/\mathbb C$, with Poisson bracket
\[
 \{f,g\}=f_u g_v-f_v g_u.
\]
Give it the quotient of the $\mathfrak m$-adic linear topology,
with discrete scalar field. The pointed subalgebra
$\mathfrak h_0=\mathfrak m^2$ has finite Lie quotients
$\mathfrak m^2/\mathfrak m^{d+1}$ for $d\geq2$.
The full algebra has finite Taylor vector-space quotients, but
these are not Lie quotients: $\{v,u^{d+1}\}=-(d+1)u^d$.

A continuous linear functional on $\mathfrak h$ kills some
$\mathfrak m^{d+1}$. Therefore
\[
 P=\mathfrak h^\vee_{\mathrm{cont}}
    =\bigoplus_{a+b>0}\mathbb C\rho_{ab},\qquad
 \rho_{ab}=u^{-a-1}v^{-b-1}\,du\,dv,
\]
with $\langle\rho_{ab},u^iv^j\rangle=\delta_{ai}\delta_{bj}$.
The full principal-part module is the local-cohomology quotient
\[
 P^{\mathrm{full}}=
 \frac{\mathcal O[u^{-1},v^{-1}]}
      {\mathcal O[u^{-1}]+\mathcal O[v^{-1}]}\,du\,dv.
\]

Its representatives have finite support in the displayed negative
monomials. The module $P^{\mathrm{full}}$ also contains
$\rho_{00}$. The subspace $P$ is the annihilator of constants
under the residue pairing. It is a coadjoint submodule, not a module over
$\mathcal O$, since $u\rho_{10}=\rho_{00}$.

The transpose of a Taylor truncation is an inclusion of finite
dual spaces. Its direct limit is $P$. In contrast,
$\prod_{a+b>0}\mathbb C\rho_{ab}$ is a larger vector space,
formed using projections in the opposite direction. It is not
the continuous dual just defined.

Fix the coadjoint convention
$\langle f\cdot\rho,g\rangle=-\langle\rho,\{f,g\}\rangle$.
Pairing against a monomial gives
\begin{equation}\label{mht:eq:full-residue-action}
 u^pv^q\cdot\rho_{ab}
   =\bigl(q(a+1)-p(b+1)\bigr)
                    \rho_{a-p+1,b-q+1}.
\end{equation}
Negative target indices give zero. The coefficient of a scalar
target also vanishes.

For $a+b\leq n$, only Hamiltonian terms of
degree at most $n+2$ can contribute, and the output degree is at
most $n+1$. This proves that a formal Hamiltonian acts on each
finite principal part through a finite Taylor window. With the
discrete linear topology on $P$, it also proves continuity of
the bilinear action.

\section{Polynomial and residual representations}

Let $\mathfrak h_{\mathrm{pol}}=\mathbb C[u,v]/\mathbb C$.
It acts on $M=\mathbb C[u,v]/\mathbb C$ by the adjoint action.
For polynomial monomial labels,
\[
 u^pv^q\cdot u^av^b=(pb-qa)u^{a+p-1}v^{b+q-1}
                         \pmod{\mathbb C}.
\]
In particular $u$ and $v$ act locally nilpotently on $M$.
This polynomial space is not preserved by every formal
Hamiltonian. For example, $f=\sum_{n\geq2}u^n$ has
$\{f,v\}=\sum_{n\geq2}n u^{n-1}$, which is not a polynomial.

\begin{theorem}\label{mht:thm:residue-simple}
The $\mathfrak h_{\mathrm{pol}}$-module $P$ is simple. If $N$ is
any module on which $u$ acts locally nilpotently, then
\[
 \operatorname{Hom}_{\mathfrak h_{\mathrm{pol}}}(P,N)=0,
 \qquad
 \operatorname{Hom}_{\mathfrak h_{\mathrm{pol}}}(N,P)=0.
\]
\end{theorem}

\begin{proof}
Take a nonzero vector in a submodule of $P$, and let $n$ be the
largest total degree of its nonzero coefficients. Choose a top
coefficient $c_{ab}\ne0$, with $a+b=n$. Acting by
$u^av^{b+1}$ kills every lower-degree term. Its output has degree
one, and its coefficient at $\rho_{10}$ is $(b+1)c_{ab}\ne0$.

The quadratic Hamiltonians act irreducibly on
$\mathbb C\rho_{10}\oplus\mathbb C\rho_{01}$: $uv$ has
eigenvalues $1,-1$, while $u^2\rho_{10}=-2\rho_{01}$ and
$v^2\rho_{01}=2\rho_{10}$. Thus the submodule contains both
degree-one vectors. Repeated action by $u$ and $v$ reaches
every $\rho_{ab}$ with $a+b>0$, with nonzero factorial
coefficients. The submodule is all of $P$.

The action of $u$ on $P$ is injective:
$u\rho_{ab}=-(b+1)\rho_{a,b+1}$. A nonzero map $P\to N$
would be injective by simplicity, contradicting local nilpotence
on its image. A map $N\to P$ sends each vector killed by a power
of $u$ to a vector killed by that power in $P$, hence to zero.
\end{proof}

Simplicity is essential for the first direction of the proof.
An injective operator on a source can acquire nilpotent action
on a quotient. Non-nilpotence alone does not rule out maps to
locally nilpotent modules.

Diagonal rescaling cannot remove the opposite index directions.
If $\rho'_{ab}=\lambda_{ab}\rho_{ab}$ with $\lambda_{ab}\ne0$, then
\[
 u\rho'_{ab}=-(b+1)\frac{\lambda_{ab}}{\lambda_{a,b+1}}
                       \rho'_{a,b+1}.
\]
Polynomial labels instead move from $b$ to $b-1$. The
distinction concerns the representation, before any choice of
factorials or deformation coefficient.

\section{Moyal coefficients and residue transposition}

Define
\[
 \mathcal P=\partial_u\otimes\partial_v
                   -\partial_v\otimes\partial_u,
 \qquad f*_h g=m\exp(\tfrac h2\mathcal P)(f\otimes g).
\]
For polynomial inputs the series is finite. Interchanging inputs
negates $\mathcal P$, so the normalized commutator is
\[
 \{f,g\}_h=
 \sum_{\substack{r\geq1\\r\ \mathrm{odd}}}
       \frac{h^{r-1}}{2^{r-1}r!}\mathcal P^r(f,g).
\]
The coefficients begin $1,1/24,1/1920,1/322560$.

For $f=u^pv^q$ and $g=u^mv^n$, write
\[
 C_r(p,q;m,n)=\sum_{j=0}^r(-1)^j\binom rj
                  (p)_{r-j}(q)_j(m)_j(n)_{r-j},
\]
where $(a)_j=a(a-1)\cdots(a-j+1)$ and $(a)_0=1$.
Then $\mathcal P^r(f,g)=C_r u^{p+m-r}v^{q+n-r}$.

Residue transposition gives
\[
 f\cdot_h\rho_{ab}=
 -\sum_{\substack{r\geq1\\r\ \mathrm{odd}}}
 \frac{h^{r-1}}{2^{r-1}r!}
 C_r(p,q;a-p+r,b-q+r)\rho_{a-p+r,b-q+r}.
\]
Terms with negative target indices vanish. The scalar target
vanishes because every bracket with the constant test function
is zero.

For linear $f$, all terms with $r\geq3$ vanish, so
the opposite translation actions above persist without quantum
correction. For $f=u^4$ and $\rho_{11}$, the only nonzero term is
\[
 u^4\cdot_h\rho_{11}=-24h^2\rho_{04}.
\]
The polynomial adjoint action on the label $uv$ instead gives
$4u^4$. These are statements about two specified modules.

\section{The algebraic Fourier lattice and its completion}

Let $R_t=\mathbb C[[t]]$, and define the algebraic Rees Weyl
algebra
\[
 D_t=R_t\langle z_1,z_2,\partial_1,\partial_2\rangle/
 ([\partial_i,z_j]=t\delta_{ij},\ [z_i,z_j]=[\partial_i,\partial_j]=0).
\]
Algebraic here means a finite sum of Weyl monomials with
coefficients in $R_t$. Its polynomial and delta modules are
\[
 O_t=R_t[z_1,z_2]=D_t/(D_t\partial_1+D_t\partial_2),
 \qquad
 \Delta_t=D_t/(D_tz_1+D_tz_2).
\]
Write $e$ for the cyclic vector of $\Delta_t$. The basis of
$\Delta_t$ is $\partial_1^a\partial_2^be$, with finite support
in $(a,b)$.

The automorphism
$\mathsf F(z_i)=\partial_i$, $\mathsf F(\partial_i)=-z_i$
preserves the Weyl relations. The map
\[
 T:O_t\longrightarrow\Delta_t,
 \qquad z_1^az_2^b\longmapsto\partial_1^a\partial_2^be
\]
satisfies $T(Pf)=\mathsf F(P)T(f)$. This specifies the
direction of the Fourier twist.

Let $L_t=\bigoplus_{a,b\geq0}R_t\rho_{ab}$, with the usual
principal-part multiplication by $z_i$ and
$\partial_i=t\partial_{z_i}$. There is an injective $D_t$-map
\begin{equation}\label{mht:eq:fourier-lattice-map}
 \iota:\Delta_t\longrightarrow L_t,
 \qquad \partial_1^a\partial_2^be\longmapsto
       (-1)^{a+b}t^{a+b}a!b!\rho_{ab}.
\end{equation}
It follows by differentiating $\rho_{00}$. The map becomes
an isomorphism after inverting $t$, because each vector has
only finitely many principal-part labels. Neither the scalar
line nor its complementary label span is a $D_t$-submodule:
$\partial_i e$ leaves the line, and
$z_i\partial_i e=-te$ leaves the complementary span.

Parameter completion changes the assertion. Put
\[
 \widehat\Delta_t=\varprojlim_r\Delta_t/t^r\Delta_t,
 \qquad \widehat L_t=\varprojlim_r L_t/t^rL_t
                         =P^{\mathrm{full}}[[t]].
\]
A vector in either completion has coefficients $b_{ab}(t)$
such that for each $r$ only finitely many labels have
$b_{ab}\notin t^rR_t$. This is the condition that each
coefficient of $t$ have finite label support. Likewise
\[
 \widehat O_t=\mathbb C[z_1,z_2][[t]]
\]
is larger than $R_t[z_1,z_2]$.

\begin{proposition}\label{mht:prop:completed-fourier-image}
The completed map $\widehat\iota$ is injective. Image membership
for a family $\sum c_{ab}(t)\rho_{ab}$ is equivalent to two
conditions. Each $c_{ab}$ is divisible by $t^{a+b}$. The coefficients
\[
 \frac{(-1)^{a+b}c_{ab}(t)}{t^{a+b}a!b!}
\]
tend to zero $t$-adically over the label set. Inverting $t$
after completion does not make this map surjective.
\end{proposition}

\begin{proof}
Formula~\eqref{mht:eq:fourier-lattice-map} acts separately
on each coefficient. Multiplication by a nonzero power of $t$
is injective in $R_t$. The stated divisibility and decay
conditions are exactly those for its coefficientwise inverse
to define a vector in $\widehat\Delta_t$.

The vector
\[
 \sigma=\sum_{n\geq0}t^n\rho_{2n,0}
\]
belongs to $\widehat L_t$. A preimage after localization would
have coefficient $t^{-n}/(2n)!$ at $\partial_1^{2n}e$.
There is no common lower bound on these powers of $t$.
Every element of $\widehat\Delta_t[t^{-1}]$ has such a lower
bound, so no preimage exists.
\end{proof}

The decay condition is independent of divisibility. The vector
$\sum_{n\geq0}n!t^n\rho_{n,0}$ satisfies divisibility, but its
formal inverse has coefficients $(-1)^n$, which do not tend
to zero. Thus even a coefficientwise inverse with no negative
powers need not belong to the completed source.

The Fourier isomorphism holds on the algebraic modules,
and $T$ extends to their parameter completions. Its completed
principal-part image is the one in the proposition. A change
of action, a localization and a completion are three separate
operations in this comparison.

\section{Support, pairing and compatible limits}

For an open interval $I$, the continuous dual of
$C_c^\infty(I)$ is the distribution space $\mathcal D'(I)$.
The continuous dual of $C^\infty(I)$ is the compactly
supported distribution space $\mathcal E'(I)$. A smooth
function acts on either by
\[
 (aT)(\varphi)=T(a\varphi),\qquad
 (\partial_tT)(\varphi)=-T(\partial_t\varphi).
\]
These formulas define multiplication by a smooth test and
transposition of differentiation. They do not define the
product of two arbitrary distributions.

Extension by zero carries compact tests from $U$ to $V$ when
$U\subset V$. Dualizing reverses this arrow. The same variance
distinction occurs between covariant Chevalley chains and
contravariant Chevalley cochains. To identify a dual cochain
object with physical observables one must construct the
required support maps, rather than infer their direction
from the coefficient formula.

A continuous Kronecker pairing cannot pair two unrestricted
coordinate products. If $e_n$ is the $n$th coordinate vector,
then $e_n\to0$ in the product topology, whereas the requested
pairing would have $B(e_n,e_n)=1$. Reciprocal changes of dual
bases also leave the contraction of the first $n$ diagonal
terms equal to $n$. A weighted Hilbert or nuclear pairing
must therefore specify its spaces and prove its convergence.

For countable surjective inverse systems $V_d,W_d$ of discrete
vector spaces, put $V=\varprojlim_dV_d$ and
$W=\varprojlim_dW_d$. A continuous map $V\to W_d$ factors
through some input quotient $V_e$. The input cutoff can depend
on the output cutoff. When the relevant quotient projections
are onto, this gives the familiar form
\[
 \operatorname{Hom}_{\mathrm{cont}}(V,W)
   =\varprojlim_d\varinjlim_e\operatorname{Hom}(V_e,W_d).
\]

This formula records continuity, not locality in a physical
coordinate. For complexes, derived inverse limits also retain
their degree-minus-one compatibility information. The strict
cycle criterion and the derived Roos complex are distinguished
in Section~\ref{mht:sec:descent-towers}.

Finally, the polar trace of Chapter~\ref{mht:ch:neumann-bv}
is a relative trace. With its boundary maps $r,\ell$, it obeys
$\tau_{\mathcal A}d=\tau_Br$ and $\tau_B\ell=0$.
On the path model with $b(0)=r(a)$ and $b(1)=\ell(c)$,
the corrected trace includes $\int_0^1\tau_B(b_{dt})$.
Indeed its value on a differential is
\[
 \tau_Br(a)+\tau_B(b(1)-b(0))=0.
\]
The endpoint term supplies closedness. A residue value alone
does not supply that term or the physical measure used in
quantum sewing.

The coefficient constructions above fix the signs and domains
of their maps. A native comparison must additionally preserve
the full field differential, support, pairing, boundary
conditions and operations. Those requirements remain at the
points where the physical maps are constructed.
